\documentclass{beamer}
\usepackage{../pomozno/prosojnice}
\setminted{
    fontsize=\scriptsize
}

\title{Najpogostejše napake v \LaTeX-u}
\date{}

\begin{document}
\frame{\titlepage}

\begin{frame}[fragile]{Matematična okolja}
  \begin{align*}
	a x^2 + b x + c &= 0 \\
	\left(x + \frac{b}{2a}\right)^2 &= \frac{b^2-4ac}{4a^2} \\
	x_{1, 2} &= \frac{-b \pm \sqrt{b^2-4ac}}{2a}
  \end{align*}

  \begin{minted}{latex}
\begin{align*}
a_{11} x_1 + a_{12} x_2 + \cdots + a_{1n} x_n &= b_1 \\
a_{21} x_1 + a_{22} x_2 + \cdots + a_{2n} x_n &= b_2 \\
&\vdots \\
a_{m1} x_1 + a_{m2} x_2 + \cdots + a_{mn} x_n &= b_m
\end{align*}
  \end{minted}
\end{frame}

\begin{frame}[fragile]{Če uporabljate \texttt{align}, uporabite tudi \texttt{\&}}
    \begin{dobroenv}
    \begin{minted}{latex}
    \begin{align*}
        a_{11} x_1 + a_{12} x_2 + \dots + a_{1n} x_n &= b_1 \\
        a_{21} x_1 + a_{22} x_2 + \dots + a_{2n} x_n &= b_2 \\
                                                     &\vdotswithin{=} \\
        a_{n1} x_1 + a_{n2} x_2 + \dots + a_{nn} x_n &= b_n,
    \end{align*}
    \end{minted}

    \begin{align*}
        a_{11} x_1 + a_{12} x_2 + \dots + a_{1n} x_n &= b_1 \\
        a_{21} x_1 + a_{22} x_2 + \dots + a_{2n} x_n &= b_2 \\
                                                     &\vdotswithin{=} \\
        a_{n1} x_1 + a_{n2} x_2 + \dots + a_{nn} x_n &= b_n,
    \end{align*}
    \end{dobroenv}
\end{frame}

\begin{frame}[fragile]{Če uporabljate \texttt{align}, uporabite tudi \texttt{\&}}
    \begin{slaboenv}
    \begin{minted}{latex}
        \begin{align*}
            a_{11} x_1 + a_{12} x_2 + \dots + a_{1n} x_n = b_1 \\
            a_{21} x_1 + a_{22} x_2 + \dots + a_{2n} x_n = b_2 \\
                                                         \vdotswithin{=} \\
            a_{n1} x_1 + a_{n2} x_2 + \dots + a_{nn} x_n = b_n,
        \end{align*}
        \end{minted}

    \begin{align*}
        a_{11} x_1 + a_{12} x_2 + \dots + a_{1n} x_n = b_1 \\
        a_{21} x_1 + a_{22} x_2 + \dots + a_{2n} x_n = b_2 \\
                                                     \vdotswithin{=} \\
        a_{n1} x_1 + a_{n2} x_2 + \dots + a_{nn} x_n = b_n,
    \end{align*}
    \end{slaboenv}
\end{frame}

\begin{frame}[fragile]{Če uporabljate \texttt{align}, uporabite tudi \texttt{\&}}
    \begin{slaboenv}
    \begin{minted}{latex}
        \begin{align*}
            a_{11} x_1 + a_{12} x_2 + \dots + a_{1n} x_n = b_1 \\
            a_{21} x_1 + a_{22} x_2 + \dots + a_{2n} x_n = b_1 + b_2 \\
                                                         \vdotswithin{=} \\
            a_{n1} x_1 + a_{n2} x_2 + \dots + a_{nn} x_n = b_1 + b_2 + \dots + b_n,
        \end{align*}
        \end{minted}

    \begin{align*}
        a_{11} x_1 + a_{12} x_2 + \dots + a_{1n} x_n = b_1 \\
        a_{21} x_1 + a_{22} x_2 + \dots + a_{2n} x_n = b_1 + b_2 \\
                                                     \vdotswithin{=} \\
        a_{n1} x_1 + a_{n2} x_2 + \dots + a_{nn} x_n = b_1 + b_2 + \dots + b_n,
    \end{align*}
    \end{slaboenv}
\end{frame}

\begin{frame}[fragile]{Če nočete oštevilčenja, uporabite \texttt{align*}}
    \begin{dobroenv}
    \begin{minted}{latex}
        \begin{align*}
            a_{11} x_1 + a_{12} x_2 + \dots + a_{1n} x_n = b_1 \\
            a_{21} x_1 + a_{22} x_2 + \dots + a_{2n} x_n = b_2 \\
                                                         \vdotswithin{=} \\
            a_{n1} x_1 + a_{n2} x_2 + \dots + a_{nn} x_n = b_n,
        \end{align*}
        \end{minted}

    \begin{align*}
        a_{11} x_1 + a_{12} x_2 + \dots + a_{1n} x_n = b_1 \\
        a_{21} x_1 + a_{22} x_2 + \dots + a_{2n} x_n = b_2 \\
                                                     \vdotswithin{=} \\
        a_{n1} x_1 + a_{n2} x_2 + \dots + a_{nn} x_n = b_n,
    \end{align*}
    \end{dobroenv}
\end{frame}

\begin{frame}[fragile]{Če nočete oštevilčenja, uporabite \texttt{align*}}
    \begin{slaboenv}
    \begin{minted}{latex}
        \begin{align}
            a_{11} x_1 + a_{12} x_2 + \dots + a_{1n} x_n = b_1 \\
            a_{21} x_1 + a_{22} x_2 + \dots + a_{2n} x_n = b_2 \\
                                                         \vdotswithin{=} \\
            a_{n1} x_1 + a_{n2} x_2 + \dots + a_{nn} x_n = b_n,
        \end{align}
        \end{minted}

    \begin{align}
        a_{11} x_1 + a_{12} x_2 + \dots + a_{1n} x_n = b_1 \\
        a_{21} x_1 + a_{22} x_2 + \dots + a_{2n} x_n = b_2 \\
                                                     \vdotswithin{=} \\
        a_{n1} x_1 + a_{n2} x_2 + \dots + a_{nn} x_n = b_n,
    \end{align}
    \end{slaboenv}
\end{frame}

\begin{frame}[fragile]{Če nočete oštevilčenja, uporabite \texttt{align*}}
    \begin{slaboenv}
    \begin{minted}{latex}
    \begin{align}
        a_{11} x_1 + a_{12} x_2 + \dots + a_{1n} x_n = b_1 \nonumber\\
        a_{21} x_1 + a_{22} x_2 + \dots + a_{2n} x_n = b_2 \nonumber\\
                                                        \vdotswithin{=}  \nonumber\\
        a_{n1} x_1 + a_{n2} x_2 + \dots + a_{nn} x_n = b_n, \nonumber
    \end{align}
    \end{minted}

    \begin{align}
        a_{11} x_1 + a_{12} x_2 + \dots + a_{1n} x_n = b_1 \nonumber\\
        a_{21} x_1 + a_{22} x_2 + \dots + a_{2n} x_n = b_2 \nonumber\\
                                                     \vdotswithin{=}  \nonumber\\
        a_{n1} x_1 + a_{n2} x_2 + \dots + a_{nn} x_n = b_n, \nonumber
    \end{align}
    \end{slaboenv}
\end{frame}

\begin{frame}[fragile]{Matematične ukaze definirajte brez \texttt{\$}}
    \begin{slaboenv}
    \begin{minted}{latex}
    \newcommand{\Rn}{$\mathbb{R}^n$}
    Kjer je matrika s koeficienti iz \Rn, $b \in$ \Rn je stolpec skalarjev,
    $x \in$ \Rn pa stolpec neznank.
    \end{minted}
    \newcommand{\Rn}{$\mathbb{R}^n$}
    Kjer je matrika s koeficienti iz \Rn, $b \in$ \Rn je stolpec skalarjev,
    $x \in$ \Rn pa stolpec neznank.
    \end{slaboenv}

    \begin{dobroenv}
    \begin{minted}{latex}
    \newcommand{\Rn}{\mathbb{R}^n}
    Kjer je matrika s koeficienti iz $\Rn$, $b \in \Rn$ je stolpec skalarjev,
    $x \in \Rn$ pa stolpec neznank.
    \end{minted}
    \newcommand{\Rn}{\mathbb{R}^n}
    Kjer je matrika s koeficienti iz $\Rn$, $b \in \Rn$ je stolpec skalarjev,
    $x \in \Rn$ pa stolpec neznank.
    \end{dobroenv}
\end{frame}

\begin{frame}[fragile]{Uporabite lahko tudi \texttt{\textbackslash ensuremath}}
    \begin{sprejemljivoenv}
    \begin{minted}{latex}
    \newcommand{\Rn}{\ensuremath{\mathbb{R}^n}}
    Kjer je matrika s koeficienti iz \Rn, $b \in \Rn$ je stolpec skalarjev,
    $x \in \Rn$ pa stolpec neznank.
    \end{minted}
    \newcommand{\Rn}{\ensuremath{\mathbb{R}^n}}
    Kjer je matrika s koeficienti iz \Rn, $b \in \Rn$ je stolpec skalarjev, $x \in \Rn$ pa stolpec neznank.
    \end{sprejemljivoenv}

    \begin{dobroenv}[Bolje]
    \begin{minted}{latex}
    \newcommand{\Rn}{\mathbb{R}^n}
    Kjer je matrika s koeficienti iz $\Rn$, $b \in \Rn$ je stolpec skalarjev,
    $x \in \Rn$ pa stolpec neznank.
    \end{minted}
    \newcommand{\Rn}{\mathbb{R}^n}
    Kjer je matrika s koeficienti iz $\Rn$, $b \in \Rn$ je stolpec skalarjev, $x \in \Rn$ pa stolpec neznank.
    \end{dobroenv}
\end{frame}

\end{document}